\documentclass[10pt,a4paper]{amsart}
\usepackage[margin=19mm]{geometry}
\usepackage{amsmath,amssymb,mathtools}
\usepackage[scr=rsfs,cal=euler]{mathalfa}
\usepackage{xfrac}
\usepackage[unicode,hidelinks,bookmarks=false]{hyperref}
\newcommand{\ie}{i.e.\ }
\newcommand{\cf}{cf.\ }
\newcommand{\resp}{respectively\ }
\newcommand{\Subsection}{\subsection}
\providecommand{\nobreakdash}{\nobreak\hbox{-}}
\setlength{\emergencystretch}{2em}
\multlinegap0pt
\usepackage{amssymb}
\usepackage{float}
\usepackage{xcolor}
\newtheorem{proposition}{Proposition}
\title{Physics-Informed Neural Networks for Computing the Morse Index of the Critical Catenoid}
\author{Miraj Samarakkody}
\date{Source: arXiv:2606.21725v1 (CC BY 4.0)}
\begin{document}
\maketitle

\noindent\emph{Source and licence.} Physics-Informed Neural Networks for Computing the Morse Index of the Critical Catenoid. Version arXiv:2606.21725v1 (CC BY 4.0), distributed by arXiv under the Creative Commons Attribution 4.0 International licence, http://creativecommons.org/licenses/by/4.0/, as recorded in the arXiv licence metadata for this exact version and shown on the abstract page https://arxiv.org/abs/2606.21725v1. Full licence notice: \texttt{licenses/arxiv-2606.21725-CC-BY-4.0.txt}.\par\smallskip

\noindent\textit{Editorial scope.} Counted unit: the article's proposition prop:rigid (source0.tex lines 357--359). The article's complete original proof of it is delivered with it (source0.tex lines 361--371, one \texttt{proof} environment of 11 lines). No mathematics is written, repaired or re-derived: the statement and the proof are the article's own lines, in the article's own notation, and no retained line is substituted or re-flowed.  Retained uncounted as the source-local context the proof needs: lines 211--213.  Nothing was omitted from any retained span, so the package carries no omission record.  No retained line contains a citation, so the wrapper carries no bibliography and no bibliography entry.  No retained line contains a figure environment, an image-inclusion command or drawing code, so no image is delivered and the package declares no asset dependency.  Editorial formatting only, in addition: an AMS \texttt{amsart} preamble that restates the article's own declarations and macros used by the retained lines, namely the environments \texttt{proposition} on separate counters, together with the standard packages those lines need.  The retained environments print in this document as Proposition 1 in order of appearance, whereas the article prints the same items with the numbering of its own full document.  The complete native source of the article as distributed in the arXiv e-print for this exact version is bundled unchanged as \texttt{sources/203290/source0.tex}.
\begin{equation}\label{eq:zeta}
	\zeta := \langle X, \nu \rangle = \frac{c}{\cosh t}\bigl(\cosh t - t\sinh t\bigr),
\end{equation}

\begin{proposition}\label{prop:rigid}
	The free-boundary condition determines $T$ uniquely as the root of $\cosh T = T\sinh T$; the radius condition then fixes $c = 1/(T\cosh T)$; and minimality fixes the Jacobi potential to $2/\cosh^2 t$, that is, $\mu = 1$. The triple $(T, c, \mu)$ admits no free parameter.
\end{proposition}

\begin{proof}
	\emph{Free-boundary condition.} The surface meets $\partial\mathbb{B}^3$ orthogonally if and only if the radial vector $X$ is tangent to $\Sigma$ along $\partial\Sigma$, equivalently $\langle X, \nu\rangle = 0$ there. From the parametrization a direct computation gives $\langle X, \nu\rangle = (c/\cosh t)(\cosh t - t\sinh t)$ as in~\eqref{eq:zeta}, so at $t = T$ the condition is $\cosh T = T\sinh T$.

	\emph{Uniqueness of $T$.} Let $g(T) = \cosh T - T\sinh T$. Then $g(0) = 1 > 0$ and $g'(T) = -T\cosh T < 0$ for $T > 0$, so $g$ is strictly decreasing and has the single root $T \approx 1.19968$. Thus $T$ is determined uniquely.

	\emph{Determination of $c$ and $\mu$.} The condition $|X| = 1$ at $t = T$ reads $c^2(\cosh^2 T + T^2) = 1$; using $\cosh T = T\sinh T$ this simplifies to $c = 1/(T\cosh T)$. Finally, the Jacobi potential is $2/\cosh^2 t$, since
	\begin{equation}
		|h|^2 \cdot (c^2\cosh^2 t) = \frac{2}{c^2\cosh^4 t}\cdot(c^2\cosh^2 t) = \frac{2}{\cosh^2 t},
	\end{equation}
	that is, $\mu = 1$ with no free parameter. Hence a family with $\mu(s) = s$ and $T(s)\neq T$ satisfies none of these conditions for $0 < s < 1$.
\end{proof}

\end{document}
